\documentclass[12pt,a4paper]{article}

\usepackage[utf8]{inputenc}
\usepackage[T1]{fontenc}
\usepackage{amsmath,amssymb,amsthm,mathtools}
\usepackage[colorlinks=true,linkcolor=blue,citecolor=red,urlcolor=blue]{hyperref}
\renewcommand*{\HyperDestNameFilter}[1]{\jobname-#1}
\usepackage{geometry}
\geometry{margin=2.5cm}
\emergencystretch=1.5em
\usepackage{enumitem}
\usepackage{booktabs}
\usepackage{graphicx}
\usepackage{needspace}
\usepackage[capitalise,noabbrev]{cleveref}

\newtheorem{theorem}{Theorem}[section]
\newtheorem{lemma}[theorem]{Lemma}
\newtheorem{proposition}[theorem]{Proposition}
\newtheorem{corollary}[theorem]{Corollary}
\newtheorem{conjecture}[theorem]{Conjecture}
\theoremstyle{definition}
\newtheorem{definition}[theorem]{Definition}
\theoremstyle{remark}
\newtheorem{remark}[theorem]{Remark}

\title{The Riemann Hypothesis: Conclusio\\[4pt]
\large Part~V of a Five-Part Series}
\author{Lukas Geiger\\
\small Bernau, Germany\\
\small ORCID: \href{https://orcid.org/0009-0005-7296-1534}{0009-0005-7296-1534}}
\date{\today}

\begin{document}
\overfullrule=0pt
\maketitle

\begin{abstract}\footnotesize
This paper concludes a five-part investigation of the Riemann Hypothesis
through spectral analysis, arithmetic thermodynamics, and the Weil
quadratic form:
\begin{description}[font=\bfseries,leftmargin=2em]
  \item[Part~I \cite{Geiger2026partI}:] Foundations and Obstructions.
  \item[Part~II \cite{Geiger2026partII}:] Even Dominance (Branch~A).
  \item[Part~III \cite{Geiger2026partIII}:] Spectral Routes
    (Branches~B and~Z).
  \item[Part~IV \cite{Geiger2026partIV}:] Cross-Route Synthesis.
  \item[Part~V (this paper):] Conclusio.
\end{description}

\noindent We present the condensed atlas: three independent conditional
reductions of RH (Even Dominance, Hadamard Strip, CCM Microcluster),
three structural identifications
($c_\Xi \approx 0.0231$, $C^* = 64\pi^2/3$ as a $\lambda=5,7$
boundary-moment benchmark with $\lambda=11$ still open,
$g_0(m) \approx 2\pi/\log\gamma_m$),
two conditional theorems (B-1, GF1), and two
isolated open proof slots. The series title remains ``From Landscape to
Atlas''; promotion to ``From Atlas to Proof'' is governed by a five-gate
audit protocol triggered only by a substantive mathematical breakthrough.

\medskip\noindent
\textbf{Series status (v3.0, 2026-05-26; criticality-lens update 2026-05-31).}
Five-part series; all three conditional reductions unchanged from v2.2.
Structural reorganisation only; no new proof claims. A 2026 working phase adds a
\emph{heuristic} criticality lens ($\mathrm{RH}\Leftrightarrow\Lambda=0$), a sharper
common-wall localization, and a computable Laguerre--P\'olya/Herglotz diagnostic
(Part~IV~\cite{Geiger2026partIV}); these re-weight priorities and close four
\emph{concrete constructions}, but make no new proof claim.
\end{abstract}

{\footnotesize
\noindent\textbf{MSC 2020:} 11M26, 11M36, 47A75\\
\textbf{Keywords:} Riemann Hypothesis, conditional reduction, atlas,
assessment, open problems, outlook
}

\clearpage


% =============================================================================
\section{Atlas Overview}
\label{sec:atlas}
% =============================================================================

\begin{center}
\renewcommand{\arraystretch}{1.3}
\resizebox{\textwidth}{!}{%
\begin{tabular}{l|l|l|l|l}
\textbf{Route} & \textbf{Sector} & \textbf{Status} & \textbf{Key result} & \textbf{Paper} \\\hline
A (Even Dom.) & Even (Part~II) & conditional & A1--A8 chain + GF1 benchmark with $C^* = 64\pi^2/3$ for $\lambda=5,7$; $\lambda=11$ open & \cite{Geiger2026partII} \\
B (Hadamard) & Spectral (Part~III) & conditional & $A_\lambda \leq 1.0063$ under log-curvature control & \cite{Geiger2026partIII} \\
Z (CCM) & Spectral (Part~III) & conditional & C2ca.1 + C2ca.2 + C2ca.5; $g_0(m) = 2\pi/\log\gamma_m$ & \cite{Geiger2026partIII, Geiger2026zookeeper} \\
Cross-Route & Combined (Part~IV) & empirical & shared Galerkin pattern; 2026 criticality lens + sharpened common wall ($\Re s=1$ continuation) & \cite{Geiger2026partIV} \\
\end{tabular}}
\end{center}


% =============================================================================
\section{Result Summary}
\label{sec:results}
% =============================================================================

\subsection{Three Structural Identifications}

\begin{enumerate}
  \item \textbf{Hadamard curvature constant (RH-neutral):}
    \[
      c_\Xi = -\frac{1}{2}\sum_\rho
              \frac{1}{(\rho-\tfrac{1}{2})^2}
            \approx 0.0231.
    \]
    Under RH, this specialises to $\sum_{k\geq 1} 1/\gamma_k^2$.
    (Part~III~\cite{Geiger2026partIII}.)

  \item \textbf{Boundary-moment constant:}
    $C^* = 64\pi^2/3 \approx 210.55$,
    identified as a candidate via
    $16\int_0^\infty t^2/\sinh^2(t/2)\,dt$; the factor~$16$ is
    conjectured to arise from the DZT2 ladder normalisation. This is a
    $\lambda=5,7$ benchmark, with $\lambda=3$ retained as a separate
    low-bandwidth guardrail and $\lambda=11$ retained as an open
    discriminator after the $N=200$ point.
    (Part~II~\cite{Geiger2026partII}.)

  \item \textbf{Hecke-normalised gap:}
    $g_0(m) \approx 2\pi/\log\gamma_m$, classical
    Montgomery--Odlyzko.
    (Part~III~\cite{Geiger2026partIII}.)
\end{enumerate}

\subsection{Two Conditional Theorems}

\begin{itemize}
  \item \textbf{Theorem B-1} (Part~III): Under uniform log-curvature
    control $H_\lambda(\delta) \leq K\delta^2$ with
    $K \leq 1.8\times 10^{-3}$:
    \[
      A_\lambda(\delta)
      \leq B_\Xi \cdot \exp(K/4) \leq 1.0063,
    \]
    where $B_\Xi = \xi(0)/\xi(\tfrac{1}{2}) = 1.00579\ldots$ is
    analytical.

  \item \textbf{Theorem GF1} (Part~II): Under eigenmode-weighted kernel
    convergence producing the effective factor~$16$:
    \[
      \mathrm{sLI}
      = \frac{64\pi^2}{3}\,\frac{\lambda}{NL} + o(\cdot).
    \]
    This is formulated for the non-exceptional converged regime; the
    $\lambda=3$ sweep and the open $\lambda=11$ discriminator prevent an
    all-$\lambda$ universal reading.
\end{itemize}


\subsection{Two Isolated Open Proof Slots}

\begin{enumerate}
  \item \textbf{B-1 slot:} Uniform log-curvature control from
    Foundation-Lock structure. Closes Branch~B.
  \item \textbf{GF1 slot:} Kernel convergence + factor-16 identification
    from DZT2 definitions. Closes Branch~A's non-exceptional
    boundary-moment benchmark, with the low-bandwidth branch kept explicit.
\end{enumerate}


% =============================================================================
\section{Localization Tables}
\label{sec:localization}
% =============================================================================

\subsection{Li Coefficient Perspective}

\begin{center}
\renewcommand{\arraystretch}{1.2}
\resizebox{\textwidth}{!}{%
\begin{tabular}{c|l|l|l}
\textbf{Step} & \textbf{Statement} & \textbf{Status} & \textbf{Method} \\\hline
S1 & Define $\xi(s)$, functional equation & proven & definition \\
S2 & Li coefficients via Bombieri--Lagarias & proven & \cite{BombieriLagarias1999} \\
S3 & Decomposition $\lambda_n = \lambda_n^{(\infty)} + \lambda_n^{(\mathrm{fin})}$ & proven & Part~I \\
S4 & Archimedean dominance: $\lambda_n^{(\infty)} \sim \frac{n}{2}\log n$ & proven & Stirling \\
\textbf{S5} & \textbf{$\lambda_n \geq 0$ for all $n \geq 1$} & \textbf{OPEN}${}^{\ast}$ & \textbf{$\Longleftrightarrow$ RH} \\
S6 & All zeros on $\Re(s) = 1/2$ & $\Leftrightarrow$ S5 & Li (1997) \\
S7 & Prime counting: $\pi(x) = \mathrm{Li}(x) + O(\sqrt{x}\log x)$ & $\Leftarrow$ S6 & von Koch
\end{tabular}%
}
\end{center}

{\footnotesize\noindent${}^{\ast}$ S5 is equivalent to RH (Li, 1997).
The present series gives only conditional reductions.\par}


\subsection{Even Dominance Perspective}

\begin{center}
\renewcommand{\arraystretch}{1.2}
\resizebox{\textwidth}{!}{%
\begin{tabular}{c|l|l|l}
\textbf{Step} & \textbf{Statement} & \textbf{Status} & \textbf{Method} \\\hline
E1 & Weil explicit formula & proven & Weil (1952) \\
E2 & Weil quadratic form $\mathrm{QW}_\lambda$ & proven & \cite{Connes2026} \\
E3 & Even/odd decomposition of $\mathrm{QW}_\lambda$ & proven & Part~II \\
E4 & Shift Parity Lemma: each prime favors even & proven & Part~II \\
E5 & Gap diagnostics at $33$ values, $\lambda \leq 1{,}300{,}000$ & evidence / conditional & Part~II \\
E5b & Euler--Maclaurin: $\rho^{\mathrm{EM}} < 0$ for $\lambda \geq 1.2\mathrm{M}$ & proven & Part~II \\
\textbf{E6} & \textbf{Cumulative step: even dominance for all $\lambda$} & \textbf{conditional reduction} & \textbf{= A6} \\
E7 & Connes' Theorem~6.1 + Hurwitz bridge $\Rightarrow$ RH & conditional on E6 & \cite{Connes2026}
\end{tabular}%
}
\end{center}


% =============================================================================
\section{Remaining Open Problems}
\label{sec:open-problems}
% =============================================================================

\begin{remark}[OP1: Cumulative Even Dominance --- reduced to the odd-sector lower bound (v2.2)]
The cumulative even dominance problem is reduced to a single spectral
question: a quantitative lower bound for $\lambda_1^-(\lambda)$
(Asymptotic Variational Gap Conjecture, Part~II).
\end{remark}

\begin{remark}[OP2 resolved: Simplicity of $\lambda_1^+$]
Certified as simple for all $33$ certificate values by interval
arithmetic. The intra-even gap grows as $\sim\!\sqrt{\lambda}$
(Part~II).
\end{remark}

\begin{conjecture}[OP3: Analytical Proof of Index Rigidity]
The negative inertia index of $K_\sigma^{\mathrm{sym}}$ is exactly~$2$
for all $N$ sufficiently large and
$\sigma \in (1, \sigma_{\max})$.
\end{conjecture}

\begin{conjecture}[OP4: Nuclear Transfer Operator for $\xi$]
Construct a separable Hilbert space $\mathcal{V}$ and a nuclear family
$s \mapsto \mathcal{L}_s$ with
$\xi(s)/\xi(0) = \det(I - \mathcal{L}_s)$, spectral radius $< 1$ for
$\Re(s) > 1/2$, and eigenvalue $1$ at the zeros.
\end{conjecture}

\begin{conjecture}[OP5: Nevanlinna Property of $\log\det_2$]
Is $\log\det_2(I - zR)$ a Nevanlinna function if and only if all zeros
lie on the real axis?
\end{conjecture}

\noindent OP2 is resolved. OP1 is reduced to a single asymptotic step.
OP3--OP5 remain structural questions.


% =============================================================================
\section{Assessment and Outlook}
\label{sec:outlook}
% =============================================================================

\subsection{What Has Been Achieved}

\begin{enumerate}
  \item A \textbf{complete structural analysis} of the gauge-theoretic
    approach (Part~I: R1--R9, K1--K4).
  \item The \textbf{Shift Parity Lemma} and \textbf{$33$
    computer-assisted finite-range diagnostics} (Part~II).
  \item The \textbf{conditional reduction A1--A8} with analytical
    backbone (M1$''$, Euler--Maclaurin, PNT transfer) and the GF1
    boundary-moment benchmark with $C^* = 64\pi^2/3$ for $\lambda=5,7$
    (Part~II).
  \item Two \textbf{independent spectral conditional reductions}
    (Branches~B and~Z, Part~III).
  \item The \textbf{cross-route synthesis} identifying shared Galerkin
    pattern and two isolated proof slots (Part~IV).
  \item \textbf{Fifteen excluded paths} (spectral + CCM) and
    \textbf{five methodical dead ends} (Parts~III--IV).
  \item A 2026 \textbf{criticality lens} ($\mathrm{RH}\Leftrightarrow\Lambda=0$,
    heuristic) re-weighting the routes toward exact-structural mechanisms; a
    sharper \textbf{common-wall} localization (the analytic continuation across
    $\Re s=1$); a computable \textbf{Laguerre--P\'olya/Herglotz diagnostic}; and
    four closures of \emph{concrete constructions} (not routes) --- all without
    changing the conditional-reduction status (Part~IV~\cite{Geiger2026partIV}).
\end{enumerate}


\subsection{Current Atlas Status}

The series title remains ``From Landscape to Atlas: Multi-Route
Cartography of an Ongoing Expedition Toward the Riemann Hypothesis.''
All three living routes are conditional reductions; no single branch has
reached conditional-proof level under the project's audit protocol.

\textbf{Promotion protocol.} The title migration from ``Atlas'' to
``From Atlas to Proof'' is governed by a five-gate audit:
(1)~single named hypothesis $H$;
(2)~every step from $H$ to RH at theorem-level;
(3)~two independent reviews including a sceptical refuter;
(4)~2--3 day cool-off;
(5)~explicit Decision Record entry.
Default: Atlas remains.


\subsection{Future Scenarios}

\begin{itemize}
  \item \textbf{One branch dies} (its open hypothesis refuted):
    documented as dead route; the other two continue.
  \item \textbf{One branch matures unconditionally}: that proof stands
    as RH; other branches retained for independent verification.
  \item \textbf{Both/all three mature}: rare, but mathematically
    valuable as multi-route verification.
  \item \textbf{None matures, but partial progress continues}: research
    programme continues; the Landscape preserves the audit. This is the
    default outlook.
\end{itemize}


\subsection{Outlook}

Concrete next steps (not time-bound):
\begin{itemize}
  \item Rigorous derivation of the two open proof slots (B-1
    log-curvature, GF1 factor-16).
  \item Resolution of the $\lambda=11$ GF1 discriminator, validation at
    higher $\lambda$, and a theorem-level separation of the $\lambda=3$
    low-bandwidth branch.
  \item $\kappa_\infty$-asymptotics for $\lambda \to \infty$.
  \item Promotion-gate audit if a slot closes.
  \item Develop the exact-structural lines (M~$=$~Bost--Connes\,$\leftrightarrow$\,Meyer;
    P-A~$=$~global adelic trace route) that the 2026 criticality lens elevates, in
    parallel with the two margin-type proof slots (Part~IV~\cite{Geiger2026partIV}).
\end{itemize}


% =============================================================================
\section{Connections to the Wider Landscape}
\label{sec:connections}
% =============================================================================

\subsection{Connes' Program}

The work in this series fits directly into Connes' spectral approach to
RH~\cite{Connes2026}. The Shift Parity Lemma (Part~II) provides
analytical evidence for the even dominance hypothesis (Theorem~6.1). The
I-1 Normal-Family reframe (Part~III) decomposes Wall~(ii) into a single
named condition~(ii-a). The CCM Microcluster Programme (Part~III)
attacks the same wall through the Fourier model.

\subsection{Nyman--Beurling--Baez-Duarte}

The Li criterion, the Weil positivity criterion, the Nyman--Beurling
criterion, and the Baez-Duarte criterion are all equivalent to RH.
Direct implications between them (without passing through RH) are known
only in two cases: Weil $\Rightarrow$ Li
(\cite{BombieriLagarias1999}) and Nyman--Beurling $\Rightarrow$
Baez-Duarte (trivial).

\subsection{The Trace-Class Obstruction as Universal Pattern}

Obstruction~K4 (trace-class failure at $\Re(s) = 1/2$, Part~I) appears
across QFT (UV divergences), statistical mechanics (thermodynamic limit),
and NCG (zeta-function regularisation). The Even Dominance approach
circumvents this by working with finite-dimensional Galerkin blocks
rather than infinite-dimensional transfer operators.


% =============================================================================
\section{Full Summary Table}
\label{sec:full-summary}
% =============================================================================

\begin{center}
\renewcommand{\arraystretch}{1.3}
\resizebox{\textwidth}{!}{%
\begin{tabular}{l|c|l}
\textbf{Contribution} & \textbf{Status} & \textbf{Paper} \\\hline
Thermodynamic formalism for $\zeta$ & established & I \\
9 unconditional structural results (R1--R9) & proven & I \\
4 excluded paths (K1--K4) & proven & I \\
Shift Parity Lemma & proven (analytic) & II \\
Gap diagnostics ($33$ values, $\lambda \leq 1{,}300{,}000$) & evidence / conditional & II \\
Leading-Mode Cancellation ($c = 2+\sqrt{2}$) & proven & II \\
Euler--Maclaurin: $\rho^{\mathrm{EM}}(13) < 0$ & proven & II \\
GF1 benchmark: $\mathrm{sLI} \sim C^*\lambda/(NL)$ for $\lambda=5,7$; $\lambda=3$ separate, $\lambda=11$ open & conditional (Thm GF1) & II \\
I-1 Normal-Family Reframe: (ii) $\to$ (ii-a) & rigorous & III \\
Foundation-Lock + converged empirics & rigorous + numerical & III \\
Theorem B-1: $A_\lambda \leq 1.0063$ (conditional) & conditional & III \\
$c_\Xi$ Hadamard curvature (RH-neutral) & analytical & III \\
CCM three C2ca-inputs & conditional reduction & III \\
6 excluded spectral + 9 excluded CCM paths & documented & III \\
Cross-route Galerkin pattern & empirical & IV \\
11 systematically explored alternatives & documented & IV \\
Hypothesis correction chain (GF1 Iter 15--27) & documented & IV \\
\textbf{Cumulative step (A6)} & \textbf{conditional reduction (v2.2)} & II \\
\textbf{Riemann Hypothesis} & \textbf{conditional reduction (v2.2)} & I--V
\end{tabular}%
}
\end{center}


\bigskip
\noindent\textbf{Acknowledgments.} Computations were performed on a
Mac~Studio (Apple M1~Max, 32~GB) and a Hetzner CCX13 cloud server
(2~dedicated vCPU, 8~GB). The interval arithmetic library mpmath was
developed by Fredrik Johansson.


\section*{AI Disclosure}

This manuscript was developed in extensive collaboration with AI language
models (Claude, Anthropic; Copilot, Microsoft/GitHub; Gemini, Google
DeepMind). AI contributed to conceptual exploration, analytical work,
literature review, writing, and critical review. The author, Lukas
Geiger, conceived the research direction, provided domain expertise, and
exercised final editorial authority. The author assumes full and sole
scholarly responsibility for all content, including any errors.


\bibliographystyle{plain}
\bibliography{fst-rh-references}

\end{document}
